\documentclass[../../../include/open-logic-section]{subfiles}

\begin{document}
\olfileid{sth}{card-arithmetic}{expotough}

\olsection[برخی ساده‌سازی‌ها]{برخی ساده‌سازی‌ها در توان‌رسانی کاردینالی}

گرچه تعریف $\canonord$ اندکی پیچیده بود، حاصل آن نتیجه‌ای سودمند دربارهٔ
جمع و ضرب کاردینالی، یعنی \olref[simp]{cardplustimesmax}، است. بااین‌حال،
توان‌رسانی ترامتناهی را نمی‌توان به همان سادگی ساده کرد. برای توضیح علت،
با نتیجه‌ای آغاز می‌کنیم که الگویی آشنا از حالت متناهی را گسترش می‌دهد
(هرچند اثبات آن در سطح بالایی از انتزاع قرار دارد):

\begin{prop}\ollabel{simplecardexpo}
$\cardexpo{\cardfont{a}}{\cardfont{b} \cardplus \cardfont{c}} =
\cardexpo{\cardfont{a}}{\cardfont{b}} \cardtimes
\cardexpo{\cardfont{a}}{\cardfont{c}}$ و
$\cardexpo{(\cardexpo{\cardfont{a}}{\cardfont{b}})}{\cardfont{c}} =
\cardexpo{\cardfont{a}}{\cardfont{b} \cardtimes \cardfont{c}}$، برای
هر سه کاردینال $\cardfont{a}, \cardfont{b}, \cardfont{c}$.
\end{prop}

\begin{proof}
برای ادعای نخست، تابعی چون $f \colon
(\cardfont{b}\disjointsum\cardfont{c}) \to \cardfont{a}$ را در نظر
بگیرید. اکنون آن را «دو پاره» کنید: برای هر
$\beta \in \cardfont{b}$ بگذارید $f_\cardfont{b}(\beta) = f(\beta, 0)$،
و برای هر $\gamma \in \cardfont{c}$ بگذارید
$f_\cardfont{c}(\gamma) = f(\gamma, 1)$. نگاشت
$f \mapsto \tuple{f_{\cardfont{b}}, f_{\cardfont{c}}}$ یک !!a{bijection}
از $\funfromto{\cardfont{b} \disjointsum \cardfont{c}}{\cardfont{a}}$ به
$(\funfromto{\cardfont{b}}{\cardfont{a}} \times
\funfromto{\cardfont{c}}{\cardfont{a}})$ است.

برای ادعای دوم، تابعی چون $f \colon \cardfont{c} \to
(\funfromto{\cardfont{b}}{\cardfont{a}})$ را در نظر بگیرید؛ پس برای هر
$\gamma \in \cardfont{c}$ تابعی چون
$f(\gamma) \colon \cardfont{b} \to \cardfont{a}$ داریم. اکنون برای هر
$\tuple{\beta, \gamma} \in \cardfont{b} \times \cardfont{c}$ تعریف کنید
$f^*(\beta, \gamma) = (f(\gamma))(\beta)$. نگاشت $f \mapsto f^*$ یک
!!a{bijection} از
$\funfromto{\cardfont{c}}{(\funfromto{\cardfont{b}}{\cardfont{a}})}$
به $\funfromto{\cardfont{b} \times \cardfont{c}}{\cardfont{a}}$ است.
\end{proof}

آنچه \emph{می‌خواهیم}، راهی آسان برای محاسبهٔ
$\cardexpo{\cardfont{a}}{\cardfont{b}}$ در هنگام سروکارداشتن با
کاردینال‌های نامتناهی است. نتیجهٔ زیر گامی مناسب در این جهت است:

\begin{prop}\ollabel{cardexpo2reduct}
اگر $2 \leq \cardfont{a} \leq \cardfont{b}$ و $\cardfont{b}$ نامتناهی
باشد، آنگاه $\cardexpo{\cardfont{a}}{\cardfont{b}} =
\cardexpo{2}{\cardfont{b}}$.
\end{prop}

\begin{proof}
\begin{align*}
	\cardexpo{2}{\cardfont{b}} &\leq
	\cardexpo{\cardfont{a}}{\cardfont{b}}\text{، زیرا $2 \leq \cardfont{a}$}\\
	&\leq \cardexpo{(\cardexpo{2}{\cardfont{a}})}{\cardfont{b}}
	\text{، بنا بر \olref[opps]{cantorcor}}\\
	&= \cardexpo{2}{\cardfont{a} \cardtimes \cardfont{b}}
	\text{، بنا بر \olref{simplecardexpo}} \\
	&= \cardexpo{2}{\cardfont{b}}
	\text{، بنا بر \olref[simp]{cardplustimesmax}}.
\end{align*}
\end{proof}

واقعاً نباید انتظار داشته باشیم که بتوان این را بیش از این ساده کرد،
زیرا بنا بر \olref[card-arithmetic][opps]{cantorcor}،
$\cardfont{b} < \cardexpo{2}{\cardfont{b}}$. بااین‌حال، این نتیجه به ما
نمی‌گوید وقتی $\cardfont{b} < \cardfont{a}$ است دربارهٔ
$\cardexpo{\cardfont{a}}{\cardfont{b}}$ چه بگوییم. البته اگر
$\cardfont{b}$ \emph{متناهی} باشد، می‌دانیم چه باید کرد.

\begin{prop}
اگر $\cardfont{a}$ نامتناهی و $0<n<\omega$ باشد، آنگاه
$\cardexpo{\cardfont{a}}{n} = \cardfont{a}$.
\end{prop}

\begin{proof}
$\cardexpo{\cardfont{a}}{n} = \cardfont{a} \cardtimes \cardfont{a}
\cardtimes \ldots \cardtimes \cardfont{a} = \cardfont{a}$، بنا بر
\olref[simp]{cardplustimesmax}.
\end{proof}
\noindent
افزون‌براین، در برخی حالت‌های دیگر نیز می‌توانیم اندازهٔ
$\cardexpo{\cardfont{a}}{\cardfont{b}}$ را مهار کنیم:

\begin{prop}
اگر $2 \leq \cardfont{b} < \cardfont{a} \leq
\cardexpo{2}{\cardfont{b}}$ و $\cardfont{b}$ نامتناهی باشد، آنگاه
$\cardexpo{\cardfont{a}}{\cardfont{b}} = \cardexpo{2}{\cardfont{b}}$.
\end{prop}

\begin{proof}
$\cardexpo{2}{\cardfont{b}}\leq \cardexpo{\cardfont{a}}{\cardfont{b}}
\leq \cardexpo{(\cardexpo{2}{\cardfont{b}})}{\cardfont{b}} =
\cardexpo{2}{\cardfont{b}\cardtimes\cardfont{b}} =
\cardexpo{2}{\cardfont{b}}$؛ استدلال همانند
\olref{cardexpo2reduct} است.
\end{proof}
\noindent
اما از این نقطه به بعد، امور بسیار ظریف‌تر می‌شوند.

\end{document}
